% Borrador aportado por Roger; precisión y formato revisados con asistencia de IA.
Durante esta práctica, mi responsabilidad fue trabajar con el reloj lógico de
Lamport, integrar su funcionamiento en cada nodo, registrar los eventos y verificar
los resultados. Aprendí que cada nodo necesita su propio reloj porque los procesos
funcionan de manera independiente. También comprendí que el servidor y el cliente
de un mismo nodo deben compartir el contador para mantener un orden coherente.
Antes no tenía tan claro cómo se relacionaban los eventos entre distintos procesos,
pero ahora entiendo mejor el concepto de causalidad.

Una de las dificultades fue diferenciar un error del programa de un fallo esperado
durante las pruebas. Cuando se detuvo el nodo C, observé que B registró el envío
fallido y que su reloj avanzó de 7 a 8. Al revisar los registros, comprobé que la
comunicación entre A y B continuaba funcionando, lo que me permitió entender mejor
el comportamiento del sistema.

En cada corrida funcional de TCP y gRPC se completaron 20 solicitudes y 20 acuses,
con 80 eventos y 40 pares, sin errores de causalidad detectados por el verificador.
Esta experiencia me permitió comprender que Lamport ordena eventos, pero no
sincroniza la hora real. Considero que una mejora sería automatizar más
verificaciones para facilitar la detección de errores.
